% sections/results.tex

\subsection{The GCSF/EPO Ratio, Not EPO Dose Alone, Controls Erythroid Commitment Probability}
\label{sec:ratio_sweep}

The central quantitative result of this work is that erythroid commitment probability in the model is governed by the \emph{ratio} of GCSF to EPO receptor occupancy, not by the absolute EPO concentration. To characterise the bifurcation structure, we ran Phase~D (7 conditions) followed by Phase~D-ext (3 fine-grain conditions): N=50 stochastic replicates each, EPO\,=\,1.0~mM fixed, GCSF\,$\in$\,\{0.5, 1.0, 1.2, 1.35, 1.5, 2.0, 2.2, 2.5, 2.8, 3.0\}~mM. The result is an abrupt, all-or-nothing switch (Table~\ref{tab:commitment_D}; Fig.~\ref{fig:p_erythroid}): P\,(ERY)\,$\geq$\,0.92 for all GCSF\,$\leq$\,2.2 (48--50/50 erythroid), switching to P\,(ERY)\,=\,0 at GCSF\,=\,2.5 (0/50 erythroid, all 50 myeloid). The commitment threshold GCSF*\,$\in$\,(2.2,\,2.5)~mM (GCSF/EPO\,$\approx$\,2.3).

\begin{table}[ht]
\caption{Phase~D and Phase~D-ext GCSF/EPO ratio sweep (N\,=\,50 per condition, EPO\,=\,1.0~mM fixed, $T=7200$~s). GATA1$_\text{nuc}$ and PU.1$_\text{nuc}$ are mean\,$\pm$\,SD across 50 replicates at $t=3600$~s. P(ERY) is the fraction of replicates with GATA1\_nuc\,$>$\,1.5\,$\times$\,PU.1\_nuc; 95\% CI uses the Wilson score. The transition between complete erythroid (P=0.96) and complete myeloid (P=0.00) commitment occurs between GCSF\,=\,2.2 and GCSF\,=\,2.5~mM (GCSF/EPO\,$\in$\,(2.2,\,2.5)).}
\label{tab:commitment_D}
\centering
\small
\begin{tabularx}{\linewidth}{@{}ccXXcc@{}}
\toprule
\textbf{GCSF (mM)} & \textbf{GCSF/EPO} & \textbf{GATA1\textsubscript{nuc} (mM)} & \textbf{PU.1\textsubscript{nuc} (mM)} & \textbf{P(ERY)} & \textbf{95\% CI} \\
\midrule
0.5 & 0.5 & $9.36 \pm 1.81$ & $0.003 \pm 0.018$ & 1.00 & [0.929, 1.00] \\
1.0 & 1.0 & $9.25 \pm 1.43$ & $0.004 \pm 0.023$ & 1.00 & [0.929, 1.00] \\
1.2 & 1.2 & $9.14 \pm 1.97$ & $0.133 \pm 0.927$ & 0.98 & [0.895, 0.996] \\
1.35 & 1.35 & $8.75 \pm 2.89$ & $0.710 \pm 2.43$ & 0.92 & [0.812, 0.968] \\
1.5 & 1.5 & $9.31 \pm 2.46$ & $0.351 \pm 1.71$ & 0.96 & [0.865, 0.989] \\
2.0 & 2.0 & $9.24 \pm 2.10$ & $0.170 \pm 1.16$ & 0.98 & [0.895, 0.996] \\
2.2 & 2.2 & $8.81 \pm 2.26$ & $0.403 \pm 1.96$ & 0.96 & [0.865, 0.989] \\
2.5 & 2.5 & $0.005 \pm 0.016$ & $8.83 \pm 1.20$ & 0.00 & [0.000, 0.071] \\
2.8 & 2.8 & $0.003 \pm 0.013$ & $8.81 \pm 1.42$ & 0.00 & [0.000, 0.071] \\
3.0 & 3.0 & $0.585 \pm 2.29$ & $8.30 \pm 2.54$ & 0.06 & [0.021, 0.162] \\
\bottomrule
\end{tabularx}
\end{table}

At GCSF\,=\,1.35~mM the commitment variance is maximal (GATA1$_\text{nuc}$ SD\,=\,2.89; PU.1$_\text{nuc}$ SD\,=\,2.43) with 46/50 erythroid replicates---a genuine mixed-fate competition zone in which individual cell outcomes are stochastic. The deep erythroid attractor (GATA1$_\text{nuc}\approx$\,9~mM, PU.1$_\text{nuc}\approx$\,0) and deep myeloid attractor (GATA1$_\text{nuc}\approx$\,0.005~mM, PU.1$_\text{nuc}\approx$\,8.8~mM at GCSF=2.5) are nearly mirror-symmetric (Fig.~\ref{fig:waddington}), consistent with a genuine bistable switch whose basin boundary lies at GCSF/EPO\,$\approx\,2.3$.

\begin{figure}[!ht]
  \centering
  \includegraphics[width=0.75\linewidth]{fig2_p_erythroid}
  \caption{\textbf{P(erythroid) vs GCSF/EPO ratio.} Erythroid commitment fraction at $t=3600$~s across ten GCSF conditions at EPO\,=\,1.0~mM fixed (N=50 per condition; Phase~D and Phase~D-ext). P(ERY) is high and roughly flat for GCSF/EPO\,$\leq$\,2.2, then switches abruptly to zero at GCSF/EPO\,=\,2.5. The commitment threshold GCSF*\,$\in$\,(2.2,\,2.5)~mM.}
  \label{fig:p_erythroid}
\end{figure}

\begin{figure}[!ht]
  \centering
  \includegraphics[width=0.75\linewidth]{fig1_waddington}
  \caption{\textbf{Erythroid and myeloid basins of attraction.} Approximate Waddington potential on the GATA1\textsubscript{nuc}$\times$PU.1\textsubscript{nuc} plane; attractor wells are estimated from EPO\,=\,1.0~mM, GCSF\,=\,1.35~mM final-state statistics (N=50), the maximum-variance condition.}
  \label{fig:waddington}
\end{figure}

\paragraph{Prior EPO-only result (Phase B) as a special case.}
The earlier Phase~B EPO sweep (0.430--0.460~mM at fixed GCSF\,=\,0.100~mM; N=50) found that P(ERY) was statistically flat across the tested EPO range (60--70\%, Kruskal--Wallis $p>0.05$). This finding is a direct consequence of the ratio-control principle: within that 7\% EPO range the GCSF/EPO ratio shifts by only $\sim$7\%, which is insufficient to cross the bifurcation boundary. Any experiment that fixes one cytokine and sweeps the other in a narrow neighbourhood of the system's operating point will appear EPO-insensitive---the apparent insensitivity is a property of the measurement window, not of the circuit.

The energy nucleotide species provided an internal consistency control for all Phase~D and Phase~D-ext simulations: ATP\,+\,ADP\,=\,33.000~mM and GTP\,+\,GDP\,=\,17.730~mM were conserved exactly across all 500 replicates ($N = 50 \times 10$ conditions), confirming numerical integrity throughout.

\subsection{A Three-Layer Cascade Architecture}

Analysis of the mean trajectory---the average of all 50 replicates in each condition---reveals that the dynamics from stimulation to settled state are organised into three functionally distinct layers whose timing and EPO-sensitivity differ sharply from one another (Fig.~\ref{fig:cascade}).

\begin{figure}[!ht]
  \centering
  \includegraphics[width=0.85\linewidth]{fig7_preemption_cascade}
  \caption{\textbf{Signal hierarchy preemption cascade.} Four $\lambda$-layer panels show sequential activation---EPO forcing ($\lambda_0$), EPOR binding ($\lambda_1$), pGATA1 commitment pulse ($\lambda_2$), and GATA1/PU.1 execution divergence ($\lambda_3$)---over $t=0$--400~s at EPO\,=\,0.449~mM.}
  \label{fig:cascade}
\end{figure}

\subsubsection*{Layer 1: Reception}

The reception layer comprises cytokine receptor binding, internalisation, and recycling at the cell surface. Its defining characteristic is a structural erythroid bias that is essentially invariant across the tested EPO range: EPO receptors are roughly four times more occupied than GCSF receptors regardless of EPO concentration, because GCSF is already saturating its receptor at the fixed dose used throughout the sweep. Doubling the EPO concentration changes EPOR occupancy by only a few percent---an effect too small to meaningfully shift the balance of downstream signalling. The reception layer therefore does not amplify EPO dose differences; it sets a coarse, constitutive erythroid bias and passes it downstream unchanged (Table~\ref{tab:layers}).

\subsubsection*{Layer 2: Decision}

The decision layer is where GATA1 and PU.1 compete for transcriptional dominance. Both genes fire at near-maximal rate from the first instant of stimulation---the system is transcriptionally primed and requires no activation lag. A persistent 4--5:1 GATA1-to-PU.1 transcription advantage provides the mean erythroid bias, but because the absolute number of mRNA molecules involved is very small, a single stochastic fluctuation---one extra PU.1 mRNA transcript---can override this average and tip the cell toward the myeloid fate. EPO dose does not modulate the transcriptional asymmetry; it is structurally set. The decision layer is therefore the locus of stochastic fate specification in this model (Fig.~\ref{fig:divergence}; Table~\ref{tab:layers}).

\begin{figure}[!ht]
  \centering
  \includegraphics[width=0.85\linewidth]{fig3_divergence}
  \caption{\textbf{Stochastic divergence (mean $\pm$1\,SD envelopes).} GATA1\textsubscript{nuc} and PU.1\textsubscript{nuc} mean\,$\pm$\,1\,SD over N=50 replicates at EPO\,=\,0.449~mM, $t=0$--3600~s; wide envelopes reflect genuine bimodality driven by the decision layer.}
  \label{fig:divergence}
\end{figure}

\subsubsection*{Layer 3: Execution}

Once GATA1 accumulates in the nucleus, the execution layer fires: GATA1 is phosphorylated by a receptor-coupled kinase, and the phosphorylated form then drives ubiquitin-mediated degradation of PU.1. This cascade activates within the first minute of stimulation and is complete within roughly two minutes. Critically, it is EPO-dose-insensitive: once triggered, it runs to the same endpoint at the same speed regardless of how much EPO is present. The execution layer does not amplify cytokine differences; it converts an upstream binary decision into a cytokine-maintained chemical record (Table~\ref{tab:layers}).

\begin{table}[ht]
\caption{Three-layer cascade summary. Functional properties of each layer as determined from the Phase B sweep (N=50 per EPO condition). GATA1:PU.1 transcription ratio refers to mean cumulative firing counts; activation and peak times are from the mean trajectory at EPO\,=\,0.449~mM.}
\label{tab:layers}
\centering
\small
\begin{tabularx}{\linewidth}{@{}lXlX@{}}
\toprule
\textbf{Layer} & \textbf{Main components} & \textbf{EPO-dose sensitivity} & \textbf{Key property} \\
\midrule
Reception & Cytokine receptor binding, internalisation, recycling & Very low ($<$6\% change across 7\% EPO range) & Structural 4:1 EPOR/GCSFR bias; set by receptor abundance \\
\addlinespace
Decision & GATA1 \& PU.1 transcription, mRNA export, translation, nuclear import & None (ratio 4--5:1 across all EPO doses) & Stochastic noise at single-digit mRNA counts determines fate \\
\addlinespace
Execution & GATA1 phosphorylation, PU.1 ubiquitin degradation & None & Activates $\sim$48~s post-stimulation; peaks $\sim$77~s; cytokine-maintained (see Phase~X) \\
\bottomrule
\end{tabularx}
\end{table}

The net result of this three-layer sequence on the mean trajectory is shown in Fig.~\ref{fig:mean_trajectory}: GATA1\textsubscript{nuc} rises continuously while PU.1\textsubscript{nuc} erodes over 3600~s, a signature that is EPO-dose-insensitive across all tested conditions.

\begin{figure}[!htb]
  \centering
  \includegraphics[width=0.85\linewidth]{fig4_mean_trajectory}
  \caption{\textbf{Mean trajectory: GATA1 rise, PU.1 erosion.} Population-mean GATA1\textsubscript{nuc} and PU.1\textsubscript{nuc} at EPO\,=\,0.449~mM over $t=0$--3600~s (N=50); the continuous GATA1 rise and PU.1 decline are the cumulative output of all three cascade layers.}
  \label{fig:mean_trajectory}
\end{figure}

\FloatBarrier
\subsection{Cytokine-Dependent Priming: GATA1/PU.1 Crossover at t\texorpdfstring{$\approx$}{~}146~s}

From the mean trajectory at EPO\,=\,0.449~mM, we identified the point at which phosphorylated nuclear GATA1 first exceeds the nuclear PU.1 protein level. This crossing occurs at approximately 146 seconds after EPO stimulation (Fig.~\ref{fig:ratchet}) and marks an important priming event: from this moment the EPOR-driven phosphorylation cascade can sustain PU.1 degradation faster than PU.1 is replenished, creating a self-reinforcing positive feedback as long as EPO receptor signalling persists. However, as demonstrated by the Phase~X cytokine-withdrawal experiment below, this state is \emph{not} an irreversible molecular record: it is maintained by ongoing EPOR signalling and collapses if that signal is removed.

\begin{figure}[!ht]
  \centering
  \includegraphics[width=0.85\linewidth]{fig5_ratchet}
  \caption{\textbf{Execution-layer priming onset ($t=0$--400\,s).} pGATA1\textsubscript{nuc} and PU.1\textsubscript{nuc} mean trajectories in the first 400~s at EPO\,=\,0.449~mM; the crossover at $t\approx146$~s marks the cytokine-dependent priming point at which GATA1 dominance is established under EPO signalling.}
  \label{fig:ratchet}
\end{figure}

Before this crossing the system is in principle reversible: a large stochastic burst of PU.1 mRNA could still redirect the cell. After the crossing, the active EPOR signalling cascade maintains GATA1 dominance; however, the system lacks an epigenetic chromatin layer and therefore remains sensitive to cytokine withdrawal (see Phase~X below). The state of the system at key timepoints under continuous EPO is summarised in Table~\ref{tab:ratchet}.

\begin{table}[ht]
\caption{Commitment state at key timepoints (mean trajectory, EPO\,=\,0.449~mM, continuous stimulation). GATA1$_\text{nuc}$ includes total unphosphorylated plus phosphorylated (pGATA1) nuclear protein; phosphorylated fraction is shown separately. PU.1 erosion reflects sustained EPOR-driven transcriptional suppression.}
\label{tab:ratchet}
\centering
\small
\begin{tabularx}{\linewidth}{@{}lcccc@{}}
\toprule
\textbf{Event} & \textbf{Time} & \textbf{GATA1\textsubscript{nuc}} & \textbf{pGATA1 fraction} & \textbf{PU.1\textsubscript{nuc}} \\
\midrule
EPO stimulation & 0~s & baseline & --- & \textasciitilde\,0.416 \\
Execution layer activated & $\sim$48~s & rising & --- & declining \\
PU.1 ubiquitin cascade peaks & $\sim$77~s & rising & --- & declining \\
GATA1/PU.1 priming crossover & $\sim$146~s & \multicolumn{2}{c}{pGATA1 = PU.1} & declining \\
Committed steady state & 3600~s & 1.04 & 59\% & 0.26 \\
\bottomrule
\end{tabularx}
\end{table}

To assess robustness of the 146~s prediction to uncertainty in the phosphorylation rate $k_{32}$, a local sensitivity analysis was performed by rescaling $k_{32}$ by factors $\alpha \in \{0.50, 0.80, 1.20, 1.50\}$. A $\pm$20\% perturbation shifts the priming-crossover time by only $\pm$3\% (143--150~s), while a $\pm$50\% perturbation shifts it by $-$40\% to $+$49\% (87--217~s). The qualitative prediction---an early cytokine-dependent priming event within the first two minutes of EPOR activation---is preserved across the entire range tested (Fig.~\ref{fig:ratchet_sensitivity}).

\begin{figure}[!htb]
  \centering
  \includegraphics[width=0.82\linewidth]{fig9_ratchet_sensitivity}
  \caption{\textbf{Priming-crossover time robustness to $k_{32}$ perturbation.}
  \textbf{A,} pGATA1$_{\text{nuc}}$ (solid) and original PU.1$_{\text{nuc}}$ (dotted) trajectories
  for $t = 0$--400~s under five $k_{32}$ scale factors $\alpha \in \{0.5, 0.8, 1.0, 1.2, 1.5\}$;
  vertical lines mark the corresponding priming-crossover times.
  \textbf{B,} Priming-crossover time vs $\alpha$; horizontal dashed line marks the directly
  measured baseline (146~s from the mean trajectory).  Even at $\pm 50\%$ perturbation
  the priming crossover occurs within 4 minutes of EPO stimulation.}
  \label{fig:ratchet_sensitivity}
\end{figure}

\subsection{Phase~X: EPO Withdrawal Reverses Erythroid Commitment}

To test whether the erythroid attractor reached under EPO represents a true irreversible state, we performed a cytokine-withdrawal experiment (Phase~X). EPO was applied at 1.0~mM for 500~s (well past the $t\approx146$~s priming crossover; GATA1$_\text{nuc}=10.91$ at $t=490$~s), then removed by setting EPO\_external\,=\,0 via a timed environment event. The trajectory was recorded to $t=3600$~s (Table~\ref{tab:phaseX}).

Following EPO withdrawal the erythroid state was not maintained. By $t=1000$~s the GATA1/PU.1 ratio had fallen from 5.97 to 3.00; by $t=2000$~s GATA1$_\text{nuc}\approx 0$ and PU.1$_\text{nuc}=1.996$, constituting a complete fate reversal to the myeloid/uncommitted attractor. The final state at $t=3600$~s was GATA1$_\text{nuc}\approx 0$, PU.1$_\text{nuc}=1.653$ (Table~\ref{tab:phaseX}).

\begin{table}[ht]
\caption{Phase~X trajectory: EPO\,=\,1.0~mM for $t=0$--500~s, then EPO\,=\,0. GATA1 dominance established at $t=490$~s collapses progressively following cytokine withdrawal.}
\label{tab:phaseX}
\centering
\small
\begin{tabularx}{\linewidth}{@{}lccc@{}}
\toprule
\textbf{Event} & \textbf{Time} & \textbf{GATA1\textsubscript{nuc}} & \textbf{PU.1\textsubscript{nuc}} \\
\midrule
EPO stimulation & 0~s & baseline & \textasciitilde\,0.385 \\
Peak erythroid (pre-withdrawal) & 490~s & 10.91 & \textasciitilde\,0 \\
EPO withdrawn & 500~s & --- & --- \\
GATA1/PU.1 ratio = 3.00 & 1000~s & declining & 1.33 \\
Full fate reversal & $\sim$2000~s & $\approx$0 & 1.996 \\
Final myeloid steady state & 3600~s & $\approx$0 & 1.653 \\
\bottomrule
\end{tabularx}
\end{table}

The mechanistic basis is direct: EPOR occupancy (EPOR\_bound) is the dominant multiplicative driver of transition T11 (GATA1 transcription) through the term $1 + 2\times\text{EPOR\_bound}/(5+\text{EPOR\_bound})$. Without EPO, EPOR\_bound decays to zero, reducing T11 to its bare self-activation rate, which is insufficient to replenish GATA1 against constitutive degradation. PU.1 then rebounds via unimpeded basal transcription.

This result demonstrates that the model's erythroid attractor is an EPOR-signalling-dependent state, not an epigenetically sealed absorbing state. True transcriptional irreversibility---the kind that produces lineage-locked cells in vivo---requires chromatin remodelling (histone methylation, accessibility changes) that deposits a memory of commitment independent of ongoing receptor occupancy. That layer is absent from the present model; its addition constitutes the highest-priority structural extension (see Section~\ref{sec:discussion}).

\subsection{Phase~E: A Minimal EPO Pulse Window Reveals Non-Monotone Commitment Stability}

To map the relationship between EPO pulse duration and post-withdrawal fate maintenance, we performed a systematic pulse-duration sweep (Phase~E). EPO was applied at 1.0~mM and removed by setting EPO\_external\,=\,0 at $t_{\mathrm{pulse}} \in \{60, 120, 300, 500\}$~s; the trajectory was then recorded to $t=3600$~s (Table~\ref{tab:phaseE}).

\begin{table}[ht]
\caption{Phase~E pulse-duration sweep. EPO\,=\,1.0~mM applied for $t_{\mathrm{pulse}}$~s then removed; GATA1$_\text{nuc}$ and PU.1$_\text{nuc}$ at $t=3600$~s. The GATA1/PU.1 threshold $>$\,1.5 is used to classify erythroid maintenance. Fate outcome is non-monotone: only the 120~s pulse maintains ERY. The ghost-attractor reversal time (last column) scales systematically with pulse length.}
\label{tab:phaseE}
\centering
\small
\begin{tabularx}{\linewidth}{@{}cccclc@{}}
\toprule
\textbf{$t_{\mathrm{pulse}}$ (s)} & \textbf{GATA1\textsubscript{nuc} (mM)} & \textbf{PU.1\textsubscript{nuc} (mM)} & \textbf{GATA1/PU.1} & \textbf{Fate} & \textbf{Reversal (s)} \\
\midrule
60  & $\approx$0.000 & 2.563 & 0.00 & Myeloid / reverted & --- (basin not entered) \\
90  & 0.002 & 3.616 & 0.00 & Myeloid / reverted & $\sim$1000 \\
\textbf{120} & \textbf{2.170} & \textbf{0.309} & \textbf{7.01} & \textbf{Erythroid (maintained)} & \textbf{---} \\
\textbf{150} & \textbf{2.704} & \textbf{0.065} & \textbf{41.4} & \textbf{Erythroid (marginal)} & \textbf{---} \\
300 & 0.019 & 2.761 & 0.01 & Myeloid / reverted & $\sim$2000 \\
500 & $\approx$0.000 & 1.653 & 0.00 & Myeloid / reverted & $\sim$2000 \\
1000 & $\approx$0.000 & 3.769 & 0.00 & Myeloid / reverted & $\sim$2500 \\
1800 & $\approx$0.000 & 0.684 & 0.00 & Myeloid / reverted & $\sim$3000 \\
\bottomrule
\end{tabularx}
\end{table}

The result is non-monotone. A 60~s pulse is insufficient to erase PU.1 (PU.1$_\text{nuc}\approx$1.8~mM at withdrawal) and the erythroid basin is never entered; PU.1 rebounds within 200~s. A 90~s pulse is a near-miss: PU.1 is 0.11~mM at withdrawal (not fully suppressed), and although GATA1$_\text{nuc}$ briefly recovers to 2.31~mM at $t=750$~s---within the target window---PU.1 is already 0.27~mM and rebuilding; the system tips to NEUT by $t\approx$1000~s. A 120~s pulse drives PU.1$_\text{nuc}$ to 0.007~mM before withdrawal, leaving GATA1 at approximately 2.0~mM; GATA1 self-activation through pGATA1 transcriptional feedback is sufficient to maintain this state without ongoing EPOR occupancy, and the final state at $t=3600$~s is stably erythroid (GATA1$_\text{nuc}=2.170$, PU.1$_\text{nuc}=0.309$, ratio~7.01). A 150~s pulse also maintains the erythroid fate (ratio~41.4 at $t=3600$~s) but by a more marginal trajectory: post-withdrawal GATA1 decays to 1.59~mM at $t=1000$~s---just above the self-activation threshold---before recovering to 2.70~mM by $t=3600$~s; PU.1 remains suppressed throughout. The lower and upper bounds of the stable ERY window are therefore bracketed: lower bound $\in$ (90,\,120)~s; upper bound $\in$ (150,\,300)~s. The confirmed ERY-maintaining interval is $t_{\mathrm{pulse}} \in [120,\,150]$~s. Pulses of 300~s, 500~s, 1000~s, and 1800~s all drive GATA1$_\text{nuc}$ to $\geq$10~mM and PU.1 to near zero before withdrawal---but all subsequently revert. The ghost-attractor reversal timescale shifts progressively later with longer pulses: $t\approx$2000~s for 300 and 500~s-pulses; $t\approx$2500~s for the 1000~s-pulse; $t\approx$3000~s for the 1800~s-pulse. This systematic scaling directly reflects the larger GATA1 overshoot produced by longer pulses, which prolongs lingering near the unstable manifold without preventing the eventual crossing. The 1000~s trajectory provides the clearest mechanistic illustration: GATA1$_\text{nuc}$ falls to 2.93~mM at $t=1100$~s and 1.48~mM at $t=1200$~s---values within the target window---yet PU.1$_\text{nuc}$ is already 0.20~mM at $t=1200$~s and continues to rebuild, confirming that the trajectory \emph{transits} the target window without arresting in it.

The mechanistic basis of the non-monotonicity is a ghost-attractor effect. A stable post-withdrawal erythroid state requires GATA1$_\text{nuc}\in(1.5,\,3)$~mM with PU.1$_\text{nuc}\lesssim$\,0.3~mM at the moment of EPO removal---the phase-space region where GATA1 self-activation dominates basal degradation without receptor input. The 120~s pulse deposits the system precisely in this window with PU.1 fully suppressed. Longer pulses drive GATA1 far above this window; the subsequent decay trajectory passes through the target window at speed while PU.1 simultaneously recovers, together ensuring that the trajectory crosses the unstable manifold rather than being captured inside the ERY basin. The critical determinant is not whether GATA1 passes through the right concentration range but whether PU.1 is sufficiently suppressed while it does so.

This result reframes the experimental prediction stated in Discussion Prediction~1: the relationship between EPO pulse duration and erythroid fate maintenance is non-monotone, with a narrow, fully bracketed window---lower bound $\in$ (90,\,120)~s; upper bound $\in$ (150,\,300)~s; confirmed ERY-maintaining interval $[120,\,150]$~s---in which stable post-withdrawal erythroid fate is achievable. Both shorter and longer EPO exposures fail, by distinct mechanisms. The complete eight-condition sweep (60, 90, 120, 150, 300, 500, 1000, 1800~s) establishes this non-monotone structure and the ghost-attractor reversal-time scaling, together constituting a double-barrelled quantitative prediction accessible to continuous live-cell GATA1/PU.1 imaging.

\subsection{The EPO* Threshold is a Population-Averaging Artefact}

The EPO* threshold---the apparent sharp transition in commitment fraction visible in bulk colony assays---emerges from the present model as a consequence of averaging across multiple stochastic replicates operating at higher copy numbers. When the model is run at 1$\times$ initial conditions (population scale), the coefficient of variation for GATA1\_nuc collapses to near zero: all replicates produce nearly identical trajectories, and the mean trajectory is a reliable predictor of individual-cell behaviour. Under these conditions, a small change in EPO concentration near the threshold switches the mean trajectory from the myeloid basin to the erythroid basin, producing the cliff-like dose-response curve that historically suggested instruction. At 0.01$\times$ initial conditions, the same circuit in the same EPO range produces P(erythroid) of 60--70\% with enormous replicate variance---the population mean masks independent draws from a bimodal distribution.

This scale dependence means that the threshold EPO* measured in bulk assays is the EPO dose that switches the majority of the population mean, which equals the EPO dose that shifts the per-cell commitment probability above 50\%. It is not a cellular switching threshold---no individual cell has a threshold; each makes a probabilistic decision influenced by noise internal to the decision layer. The EPO* is therefore a statistical property of the measurement procedure and the population copy number, not a property of the molecular circuit. This interpretation is consistent with, and provides a mechanistic explanation for, the failure of multiple groups to demonstrate EPO*-like all-or-none switching in single-cell assays~\cite{hoppe2016,olsson2016}.

\subsection{Validation Against Hoppe et\,al.\ 2016}

Hoppe et\,al.~\cite{hoppe2016} used dual-reporter live-cell imaging to track GATA1 and PU.1 protein dynamics in single myeloid progenitors and found that individual commitment decisions are stochastic with high cell-to-cell variability; population-level colony assays misrepresent this as a deterministic EPO* switch. To validate our scale-dependent noise predictions quantitatively, we compare three outputs of our two simulation runs against the parameters reported by Hoppe et\,al.\ (Fig.~\ref{fig:hoppe_validation}).

\textbf{(A) Commitment fraction.}
At $1\times$ initial conditions (population scale, $N=20$), every replicate commits to the erythroid attractor across all EPO concentrations tested (P(erythroid)\,=\,100\%, CV\,$\approx$\,0), consistent with the deterministic-like behaviour seen in bulk colony assays. At $0.01\times$ initial conditions (single-cell scale, $N=50$), the same EPO range (0.430--0.460~mM) yields P(erythroid)\,=\,60--70\% (Table~\ref{tab:commitment_D}), with the fraction statistically flat across EPO concentrations (Kruskal--Wallis $p>0.05$). This plateau at sub-maximal commitment, without a dose-dependent increase, matches the wide commitment distributions reported by Hoppe et\,al.\ under physiological EPO.

\textbf{(B) Cell-to-cell variability.}
The coefficient of variation (CV) of GATA1\textsubscript{nuc} at $t=3600$~s in the v6 single-cell run ($N=50$, EPO\,=\,1.0\,mM, GCSF\,=\,1.0\,mM) is 0.29, compared to CV\,$\approx$\,0 in the population-scale run. Hoppe et\,al.\ reported GATA1 expression CV values in the range 0.35--0.45 from single-cell time-lapse data~\cite{hoppe2016}. The v6 CV (0.29) falls slightly below this empirical window but is an order of magnitude above the population-scale value, confirming that the 0.01$\times$ copy-number regime places the model in the qualitatively correct high-noise regime for single haematopoietic progenitors. The mild quantitative shortfall likely reflects the choice of EPO\,=\,1.0\,mM (deep erythroid basin, P(ERY)\,=\,1.00), where intra-basin stochastic spread is predominant; at EPO near the bifurcation point, fate-branching uncertainty between replicates would increase the apparent CV toward the Hoppe range.

\textbf{(C) Commitment timing.}
The mean-trajectory GATA1/PU.1 crossing time (the point at which the population-mean GATA1 first exceeds mean PU.1) is $t\approx0.4$~s in the population-scale run---essentially immediate given the 1~mM scale. In the single-cell run this crossing is reached at $t=3.6$--8.3~s in mean concentration terms, but individual stochastic cells take far longer: the priming crossover ($t\approx146$~s) marks the point at which cytokine-dependent GATA1 dominance is established---a cytokine-maintained state, not an irreversible record (see Phase~X). The delay between signal (EPO binding) and priming establishment is thus an emergent property of molecular-count noise, not of kinetic rate constants, consistent with the variable commitment latencies observed by Hoppe et\,al.\ in live imaging.

Together, these three comparisons show that the scale-dependent noise architecture of the model---a deterministic macro-behaviour collapsing into broad stochastic commitment at single-cell copy numbers---is qualitatively consistent with Hoppe et\,al.\ 2016. The v6 model reproduces the correct commitment fraction pattern (P(ERY)\,$\approx$\,1.0 at saturating EPO) and places single-cell variability in the high-noise regime (CV\,=\,0.29 vs.\ population CV\,$\approx$\,0), though the absolute CV at the deep-erythroid operating point falls slightly short of the empirical window (0.35--0.45).

\begin{figure}[!htb]
  \centering
  \includegraphics[width=\linewidth]{fig8_hoppe_validation}
  \caption{\textbf{Validation against Hoppe et\,al.\ 2016.}
  \textbf{A,} Erythroid commitment fraction at $t=3600$~s.  Population-scale run
  ($1\times$ ICs, $N=20$) yields 100\% erythroid at all EPO doses.  Single-cell run
  ($0.01\times$ ICs, $N=50$) yields 60--70\% (Table~\ref{tab:commitment_D}; shaded
  band: Hoppe et\,al.\ reference range~\cite{hoppe2016}).
  \textbf{B,} Coefficient of variation (CV) of GATA1\textsubscript{nuc} at $t=3600$~s.
  Population-scale CV\,$\approx$\,0; single-cell CV\,=\,0.29 (v6, EPO\,=\,1.0\,mM)
  (ochre band: Hoppe 2016 empirical range 0.35--0.45).
  \textbf{C,} Mean-trajectory GATA1/PU.1 crossing time.  Population scale:
  $\approx$0.4~s (deterministic, near-instantaneous); single-cell scale: 3.6--8.3~s
  (noise-delayed mean crossing).}
  \label{fig:hoppe_validation}
\end{figure}

\FloatBarrier
\subsection{Nuclear pH Modulates Erythroid Commitment Probability via a pH-Hill Mechanism}

The Phase B sweep held nuclear pH constant at 7.5 throughout. Phase G-v3 extended the analysis to a factorial design spanning six EPO concentrations (0.436--0.446~mM) and three nuclear pH values (7.0, 7.5, 8.0) at GCSF\,=\,1.1~mM, using N\,=\,50 stochastic replicates per condition (18 conditions, 900 total replicates; run \texttt{run\_20260312\_233923}). The design was motivated by the pH-Hill inhibition term governing GATA1 and PU.1 transcription (model transitions T11, T12):

\begin{equation}
K_{\mathrm{inh}}(\mathrm{pH}) = 8.0 \times 10^{\,0.5(\mathrm{pH} - 7.5)} \ \mathrm{mM}
\label{eq:pHhill}
\end{equation}

\noindent At pH\,=\,7.0, $K_{\mathrm{inh}} = 4.499$~mM; at pH\,=\,7.5 (Phase B baseline), $K_{\mathrm{inh}} = 8.000$~mM; at pH\,=\,8.0, $K_{\mathrm{inh}} = 14.226$~mM. Because $K_{\mathrm{inh}}$ is an inhibition half-maximum constant, a lower value strengthens cross-inhibition in both directions. Near the bifurcation point---where PU.1 dominates---stronger PU.1-mediated inhibition of GATA1 transcription raises the EPO commitment barrier; once committed---where GATA1 dominates---stronger GATA1-mediated silencing of PU.1 deepens the erythroid attractor. The two effects operate in opposite practical directions and are addressed separately below.

Aggregated P(ERY) was monotonically increasing from pH\,=\,7.0 to pH\,=\,8.0 (Table~\ref{tab:phasegv3}), consistent with the prediction that alkaline nuclear pH modestly favours commitment near the bifurcation. However, the absolute effect was small: alkaline pH raised P(ERY) by approximately a factor of 2 across the full physiological pH range (2.0\% at pH\,=\,7.0 versus 4.0\% at pH\,=\,8.0), and standard errors overlapped at all pH pairs. The EPO commitment step was sharp and pH-independent: all 150 replicates at EPO\,=\,0.436~mM failed to commit across all three pH values, while a non-zero P(ERY) appeared at EPO\,=\,0.438~mM, placing EPO* $\approx$\,0.437~mM regardless of pH---consistent with the Phase D ratio sweep (see Section~\ref{sec:ratio_sweep}).

\begin{table}[ht]
\caption{Phase G-v3 aggregate commitment fractions by pH (run \texttt{run\_20260312\_233923}). P(ERY) is the fraction of committed replicates summed over all six EPO doses (N\,=\,50 each; total 300 replicates per pH). SEs are $\sqrt{p(1-p)/300}$.}
\label{tab:phasegv3}
\centering
\small
\begin{tabularx}{\linewidth}{@{}ccc@{}}
\toprule
\textbf{pH} & \textbf{P(ERY)} & \textbf{$\pm$\,1 SE} \\
\midrule
7.0 & 0.020 & 0.008 \\
7.5 & 0.030 & 0.010 \\
8.0 & 0.040 & 0.011 \\
\bottomrule
\end{tabularx}
\end{table}

\subsubsection*{Attractor Depth is pH-Dependent and Reversed from Earlier Predictions}

Phase H isolated the effect of nuclear pH on committed cells by running EPO\,=\,1.0~mM (far above EPO*) with GCSF\,$\approx$\,0, N\,=\,50 per pH value (150 replicates total; run \texttt{run\_20260313\_075757}). P(ERY)\,=\,100\% at all three pH values, enabling a clean within-committed comparison of attractor depth.

\begin{table}[ht]
\caption{Phase H: residual nuclear PU.1 in committed cells at EPO\,=\,1.0~mM, GCSF\,$\approx$\,0 (N\,=\,50 per pH; run \texttt{run\_20260313\_075757}). Smaller PU.1$_\text{nuc}$ indicates a deeper erythroid attractor. Medians were identically 0 for all pH values; the mean and maximum capture rare outlier replicates with measurable residual PU.1.}
\label{tab:phaseh}
\centering
\small
\begin{tabularx}{\linewidth}{@{}cXXX@{}}
\toprule
\textbf{pH} & \textbf{PU.1$_\text{nuc}$ mean (mM)} & \textbf{$\pm$\,SD} & \textbf{Max (mM)} \\
\midrule
7.0 & 0.000092 & 0.000649 & 0.0046 \\
7.5 & 0.003873 & 0.021241 & 0.143  \\
8.0 & 0.006114 & 0.030238 & 0.156  \\
\bottomrule
\end{tabularx}
\end{table}

Residual PU.1$_\text{nuc}$ increased monotonically with alkaline pH (Table~\ref{tab:phaseh}): acidic conditions produced the \emph{deepest} erythroid attractor, not the shallowest. This is mechanistically consistent with the pH-Hill $K_{\mathrm{inh}}$ term: in the committed state GATA1 dominates, so the physiologically relevant interaction is GATA1 inhibiting PU.1; at acidic pH, the smaller $K_{\mathrm{inh}}$ makes this suppression more complete. In absolute terms the differences are minor---PU.1$_\text{nuc}$ at pH\,=\,8.0 is $\approx$\,0.006~mM against $\approx$\,9~mM for GATA1$_\text{nuc}$, a ratio exceeding 1000:1 even in the most alkaline condition---so all three pH values are functionally equivalent once commitment is reached; the pH effect on attractor depth is a second-order modulation.

Phases G-v3 and H together reveal two separable regimes governed by the same $K_{\mathrm{inh}}(\mathrm{pH})$ parameter: near the bifurcation, alkaline pH modestly increases the committed fraction ($\approx$\,2$\times$ over the physiological pH range); in fully committed cells, acidic pH produces marginally more complete PU.1 silencing. Nuclear pH therefore acts as a weak, bidirectional modulator of erythroid commitment rather than a unidirectional rheostat, and its effect on commitment probability is dwarfed by the GCSF/EPO ratio (which controls P(ERY) over a $\sim$40-fold range with sharp threshold behaviour; see Section~\ref{sec:ratio_sweep}).
